% ================= 4. DYNAMIC MODELING & CONTROL =================
\section{Dynamic Modeling \& Control}
\label{sec:modeling}


The full Cubli is underactuated and simultaneously unstable about all three
axes when balanced on a corner (Section~\ref{sec:objectives}). The modeling,
identification, and control workflow is therefore first validated on a
two-dimensional (2D) single-face prototype with one rotational degree of
freedom, before being generalized to the 3D cube in
Section~\ref{sec:3dmodel}.

\subsection{2D Reaction-Wheel Pendulum}
\label{sec:2dmodel}

The 2D prototype is a 3D-printed plate sized to one cube face, with a
motor-driven momentum wheel at its center and a braking mechanism at one
corner. The plate pivots on a bearing at that corner, constraining motion to a
single rotational degree of freedom in a vertical plane.

% Suggested figure: free-body diagram of the 2D pendulum labeling
% theta_b, theta_w, l, l_b, and the pivot.

\subsubsection{Equations of Motion}
\label{sec:2deom}

The nonlinear equations of motion are
\begin{align}
\ddot{\theta}_b &= \frac{(m_b l_b + m_w l)\, g \sin\theta_b - T_m - C_b \dot{\theta}_b + C_w \dot{\theta}_w}{I_b + m_w l^2}, \\
\ddot{\theta}_w &= \frac{(I_b + I_w + m_w l^2)(T_m - C_w \dot{\theta}_w)}{I_w (I_b + m_w l^2)} \notag \\
&\quad - \frac{(m_b l_b + m_w l)\, g \sin\theta_b - C_b \dot{\theta}_b}{I_b + m_w l^2},
\end{align}
with motor current-torque relation $T_m = K_m u$. Symbols are defined in
Table~\ref{tab:2d_vars} and used throughout this section.

\begin{table}[h]
\centering
\caption{2D pendulum model variables.}
\label{tab:2d_vars}
\begin{tabular}{lll}
\toprule
Symbol & Description & Unit \\
\midrule
$\theta_b$ & Body tilt angle & rad \\
$\theta_w$ & Wheel rotation relative to body & rad \\
$m_b, m_w$ & Body, wheel mass & kg \\
$I_b, I_w$ & Body (about pivot), wheel (about motor axis) inertia & kg\,m$^2$ \\
$l$ & Motor axis to pivot distance & m \\
$l_b$ & Body CoM to pivot distance & m \\
$g$ & Gravitational acceleration & m/s$^2$ \\
$T_m$ & Motor torque & N\,m \\
$C_b, C_w$ & Body, wheel damping coefficient & kg\,m$^2$/s \\
$K_m$ & Motor torque constant & N\,m/A \\
$u$ & Motor current input & A \\
\bottomrule
\end{tabular}
\end{table}

\subsubsection{Equilibria \& Linearization}
\label{sec:2dequilibria}

The relevant equilibrium is the upright position
$(\theta_b,\dot\theta_b,\dot\theta_w)=(0,0,0)$; its linearization is presented
with the control design in Section~\ref{sec:control2d}.
\newpage
\subsection{Parameter Identification \& the 2D Testbed}
\label{sec:paramid}

CAD mass-properties simulation, accounting for print material, infill, and
mounted hardware (driver, encoder, MCU, gyroscope/accelerometer, wiring),
gives $m_b, m_w, I_b, I_w$ directly, while $l$ and $l_b$ follow from the
plate's design geometry. The damping coefficients $C_b$ and $C_w$ cannot be
obtained from CAD, since they depend on bearing friction, wiring drag, and
motor cogging; nominal values will be estimated for the initial control
design, with dedicated testing planned once the prototype is fabricated.


\begin{table}[h]
\centering
\caption{Parameter determination method (symbols per Table~\ref{tab:2d_vars}).}
\label{tab:params}
\begin{tabular}{ll}
\toprule
Symbol & Determination method \\
\midrule
$l$   & Calculated during design \\
$l_b$ & Calculated during design \\
$m_b$ & CAD mass properties \\
$m_w$ & CAD mass properties \\
$I_b$ & CAD mass properties \\
$I_w$ & CAD mass properties \\
$C_b$ & Estimated (future testing planned with prototype) \\
$C_w$ & Estimated (future testing planned with prototype) \\
$K_m$ & Motor datasheet \\
\bottomrule
\end{tabular}
\end{table}

\subsection{Jump-Up Feasibility Analysis}
\label{sec:jumpup}

Given CAD-derived $I_b, I_w, m_b, l_b, m_w, l$, the wheel speed $\omega_w$
required for jump-up is estimated before fabrication, assuming a perfectly
inelastic wheel-body collision. Conservation of angular momentum at impact and
of energy from impact to edge-standing yield
\begin{equation}
\omega_w^2 = (2-\sqrt2)\,\frac{I_w+I_b+m_wl^2}{I_w^2}\,(m_bl_b+m_wl)\,g.
\end{equation}
This value is checked against the motor's achievable speed at the available
bus voltage and used to size wheel inertia (radius, rim mass) within the
motor's operating range without excessively reducing the balancing recovery
angle. Post-fabrication, $\omega_w$ is refined experimentally to correct for
deviations from the ideal collision assumption.

The braking mechanism uses a servo-driven barrier that strikes a bolt head on
the wheel to stop it abruptly. Servo response time bounds the validity of the
``instantaneous stop'' assumption; a slower stop yields a lower effective
post-impact body velocity than predicted, compensated through the same tuning
process.
\subsubsection{Torque-Limited Wheel-Inertia Target}
\label{sec:jumpup_target}

The relation above assumes an instantaneous, lossless momentum transfer. For
preliminary sizing, the achievable motor speed is treated as fixed and the
required per-wheel inertia is derived from a momentum budget that additionally
accounts for finite brake torque. This gives an ideal per-wheel angular
momentum
\begin{equation}
h_{w,\text{ideal}} = \sqrt{\eta}\;\frac{\sqrt{2 g\,\lambda\,\kappa}}{n}\;M L^{3/2},
\label{eq:hw_ideal}
\end{equation}
where $M, L, g$ are cube mass, edge length, and gravity, $\eta$ is an energy
margin, and $\kappa,\lambda,n$ are dimensionless geometric factors set by the
contact case (Table~\ref{tab:contact_factors}). 
\newpage
The corner case is more
demanding and governs the design.

\begin{table}[h]
\centering
\caption{Contact-case geometric factors.}
\label{tab:contact_factors}
\begin{tabular}{lccc}
\toprule
Case & $\kappa$ (pivot inertia) & $\lambda$ (CoM rise) & $n$ (momentum transfer) \\
\midrule
Edge & $2/3$ & $1/\sqrt2 - 1/2$ & $1$ \\
Corner (governs design) & $11/12$ & $\sqrt3/2 - 1/2$ & $\sqrt2$ \\
\bottomrule
\end{tabular}
\end{table}

A finite brake torque $\tau_b$ cannot deliver the impulse instantaneously and
must exceed the gravitational tip-over torque $\tau_g = MgL/2$. This loss is
captured by a brake-amplification factor
\begin{equation}
\beta = \frac{\tau_b}{\tau_b - \tau_g},
\qquad
h_w = \beta\,h_{w,\text{ideal}},
\qquad
I_{w,\text{target}} = \frac{h_w}{\omega_{\max}},
\label{eq:Iw_target}
\end{equation}
where $\omega_{\max}$ is the maximum practical wheel speed.
Equations~\eqref{eq:hw_ideal} and \eqref{eq:Iw_target} are the sizing relations
used at step~4 of the closure procedure of Section~\ref{sec:massbudget}: given a
candidate cube mass and edge length they return the inertia the wheel of
Section~\ref{sec:rw_sizing} must meet. Their inputs are recorded in
Table~\ref{tab:jumpup_budget} and are not fixed here---$M$ and $L$ are outputs
of the packaging closure, not assumptions preceding it.

Two features of these relations govern how the iteration behaves. First,
$h_w \propto M L^{3/2}$, so growing the cube to accommodate a larger wheel also
raises the requirement; convergence depends on the achievable inertia growing
faster ($\sim D_w^2$) than the requirement does. Second, $\beta$ diverges as
$\tau_b \to \tau_g$: a brake whose torque only marginally exceeds the
gravitational tip-over torque cannot deliver the required impulse at any wheel
speed. The brake torque must therefore be checked against $\tau_g = MgL/2$ on
every iteration, and a comfortable ratio $\tau_b/\tau_g$ maintained rather than
a merely positive margin.

\begin{table}[h]
\centering
\caption{Torque-limited wheel-inertia target (corner case). Inputs are fixed by
the closure procedure of Section~\ref{sec:massbudget}; the table is populated on
each iteration. $\tau_b$ is an assumed value pending measurement, justified
below and equivalent to a \qty{38}{\milli\second} stop at the present $h_w$.}
\label{tab:jumpup_budget}
\begin{tabular}{llr}
\toprule
Symbol & Description & Value \\
\midrule
$M$                & Cube mass                    & \textcolor{red}{TBD} \\
$L$                & Cube edge length             & \textcolor{red}{TBD} \\
$\eta$             & Energy margin                & \textcolor{red}{TBD} \\
$\tau_b$           & Brake torque per wheel       & \qty{5}{\newton\metre} (assumed; see below) \\
$\tau_g = MgL/2$   & Gravitational tip-over floor & \textcolor{red}{TBD} \\
$\beta$            & Brake-amplification factor   & \textcolor{red}{TBD} \\
$\omega_{\max}$    & Max wheel speed (design)     & \textcolor{red}{TBD} \\
$h_w$              & Required angular momentum    & \textcolor{red}{TBD} \\
$I_{w,\text{target}}$ & \textbf{Target wheel inertia} & \textcolor{red}{\textbf{TBD}} \\
\bottomrule
\end{tabular}
\end{table}

\paragraph{Choice of $\omega_{\max}$.} The maximum wheel speed is a design
choice rather than a datasheet figure, and it enters the sizing directly
through $I_{w,\text{target}} = h_w/\omega_{\max}$: a higher assumed speed buys a
proportionally smaller wheel. It is therefore the natural relief valve when the
envelope caps the wheel below the requirement (step~5,
Table~\ref{tab:sizing_procedure}), but it must be set from what the drive can
sustain in service, not from the no-load figure $K_v V_{\text{bus}}$. Three
effects separate the two: winding resistance, which costs speed under the torque
required to spin the wheel up; the reduced phase voltage available under
sinusoidal field-oriented modulation; and the fall of pack voltage through the
discharge curve, which must be taken at its lower end if the maneuver is to
remain available on a partly-depleted pack rather than only on a fresh one. The
design value is set below the ceiling those effects imply, leaving margin for
motor tolerance, bearing and aerodynamic drag, and the residual uncertainty in
$M$. The resulting back-EMF should remain a comfortable fraction of the usable
bus voltage, and the electrical frequency $\omega_{\max}p/2\pi$ at $p$ pole
pairs must sit inside the driver's commutation capability---both checks to be
confirmed once the motor and pack are selected.

\paragraph{Choice of $\tau_b$, and what it is not.} The brake torque requires
more care than the other inputs, because the mechanism does not generate it in
the way the symbol suggests. The brake is a servo-driven barrier struck by a
bolt head on the wheel: the servo only \emph{holds} the barrier in the wheel's
path, and the wheel's own momentum does the work of the collision. The servo's
stall torque is therefore \emph{not} $\tau_b$ and need not approach it---the
TowerPro MG92B of Table~\ref{tab:bom-supplied} is rated near
\qty{0.29}{\newton\metre}, roughly one seventeenth of the value used in sizing,
and this is not an inconsistency. It must resist the barrier being deflected,
not decelerate the wheel.

Because the stop is a collision rather than a friction clutch, $\tau_b$ is not
an independent property of any component and cannot be read from a datasheet.
It is fixed entirely by how quickly the angular momentum is destroyed,
\begin{equation}
\tau_b = \frac{h_w}{\Delta t},
\label{eq:tau_b_from_dt}
\end{equation}
so quoting a brake torque is equivalent to asserting a stopping time. At the
present operating point ($h_w \approx \qty{0.19}{\kilo\gram\metre\squared\per
\second}$) the design value $\tau_b = \qty{5}{\newton\metre}$ corresponds to
$\Delta t \approx \qty{38}{\milli\second}$. Stated that way the assumption is
open to inspection: \qty{38}{\milli\second} is characteristic of servo travel
rather than of a rigid impact, and a bolt head striking a printed barrier
should arrest the wheel in a few milliseconds, implying a substantially larger
$\tau_b$. The design value is thus expected to be conservative for
\emph{sizing} purposes.

That conservatism is deliberate, and it is asymmetric. Since
$\beta = \tau_b/(\tau_b-\tau_g)$ falls towards unity as $\tau_b$ grows,
underestimating the brake torque inflates the inertia requirement---the safe
direction---while overestimating it leaves the wheel undersized. With the
present $\tau_g$, raising $\tau_b$ from \qtyrange{5}{10}{\newton\metre} relaxes
$I_{w,\text{target}}$ by only about \qty{12}{\percent}, whereas lowering it to
\qty{2}{\newton\metre} inflates the requirement by roughly \qty{65}{\percent}
and to \qty{1.5}{\newton\metre} by over \qty{150}{\percent}. The penalty for
being wrong is concentrated almost entirely on the low side, which is why a
mid-range value is adopted in preference to the higher figure a short impact
would justify.

One consequence must be carried into the structural design rather than left
implicit: $\tau_b$ is also a \emph{load}. The sizing value is chosen for the
dynamics, and the same number should not be used to size the barrier, the bolt,
the servo bracket or the wheel spokes, since a genuine few-millisecond stop
would put an order of magnitude more torque through that load path. Until
measurement settles the question, the structure is to be checked against a
short-$\Delta t$ case and the dynamics against
$\tau_b = \qty{5}{\newton\metre}$. The measurement that closes this is the
spin-down test of Section~\ref{sec:test_jumpup}, which should report both the
impulse-equivalent mean $\bar\tau_b = h_w/\Delta t$, which is the sizing
quantity, and the peak torque, which is the structural one; the two may differ
several-fold, and the logging rate must be fast enough to resolve an event of a
few milliseconds without smoothing the peak away.

\subsection{Control Design on the 2D Model}
\label{sec:control2d}

\subsubsection{LQR Baseline Control}

Linearizing about the upright equilibrium $(\theta_b,\dot\theta_b,\dot\theta_w)=(0,0,0)$ gives
\begin{equation}
\dot x(t) = Ax(t) + Bu(t), \qquad x := (\theta_b,\dot\theta_b,\dot\theta_w),
\end{equation}
with $A, B$ built from the parameters in Table~\ref{tab:2d_vars}.
Zero-order-hold discretization at a rate set by the open-loop unstable pole,
with a safety margin, gives $x[k+1] = A_dx[k]+B_du[k]$. A discrete LQR gain
$K_d=(K_{d1},K_{d2},K_{d3})$ minimizes
\begin{equation}
J(u) = \sum_k x^{\mathrm T}[k]Qx[k] + u^{\mathrm T}[k]Ru[k],
\end{equation}
giving $u[k]=-K_d(\hat\theta_b[k],\hat{\dot\theta}_b[k],\hat{\dot\theta}_w[k])$.

The tilt estimate $\hat\theta_b$ is obtained from a gyroscope and
accelerometer mounted at, or as close as mechanically feasible to, the pivot
point rather than at the plate's center. This placement is necessary because
an accelerometer at radius $r$ from the pivot measures gravity plus the
tangential ($r\ddot\theta_b$) and centripetal ($r\dot\theta_b^2$) acceleration
components induced by the plate's rotation, which corrupt the naive estimate
$\tan^{-1}(-a_x/a_y)$ during fast motion (e.g., jump-up or aggressive
recovery); placing the sensor at $r\approx0$ makes both terms negligible.
Tilt is estimated with a complementary filter combining gyro integration
(drift-prone, unaffected by the terms above) with accelerometer inclination
(drift-free, but unreliable during shocks):
\begin{equation}
\hat\theta_b[k] = (1-\gamma)\big(\hat\theta_b[k-1] + \dot\theta_{b,\text{gyro}}[k]\,\Delta t\big) + \gamma\,\theta_{b,\text{accel}}[k],
\end{equation}
where $\theta_{b,\text{accel}}[k] = \tan^{-1}(-a_x[k]/a_y[k])$ is the
instantaneous accelerometer-only estimate and $\gamma\in(0,1)$ is a small
weight favoring the gyro term between corrections.

A constant tilt-sensor offset $d$, including residual bias from imperfect
pivot placement or sensor misalignment, would otherwise appear as a nonzero
steady-state wheel speed rather than a corrected angle. A low-pass filter
state
\begin{equation}
x_f[k+1] = (1-\alpha)x_f[k] + \alpha\hat\theta_b[k]
\end{equation}
is subtracted from the tilt feedback,
\begin{equation}
u[k] = -K_d(\hat\theta_b[k]-x_f[k],\, \hat{\dot\theta}_b[k],\, \hat{\dot\theta}_w[k]),
\end{equation}
so $d$ is absorbed by $x_f$ at steady state, requiring no separate calibration
stage. This fixed-gain design is the baseline balancing controller;
Table~\ref{tab:control2d_vars} summarizes the control variables.

\begin{table}[h]
\centering
\caption{Control-design variables (2D and corner-balancing 3D model).}
\label{tab:control2d_vars}
\begin{tabular}{ll}
\toprule
Symbol & Description \\
\midrule
$x$ & State vector \\
$A, B$ & Continuous-time state/input matrices (linearized model) \\
$A_d, B_d$ & Discrete-time (ZOH) counterparts of $A, B$ \\
$K_d$ & LQR feedback gain \\
$Q, R$ & LQR state/input weighting matrices \\
$J(u)$ & LQR cost function \\
$\hat\theta_b, \hat{\dot\theta}_b$ & Estimated tilt angle, rate (gyroscope/accelerometer) \\
$\hat{\dot\theta}_w$ & Estimated wheel rate (encoder) \\
$a_x, a_y$ & Raw accelerometer readings (tangential, radial axes) \\
$\theta_{b,\text{accel}}$ & Instantaneous accelerometer-only tilt estimate \\
$\dot\theta_{b,\text{gyro}}$ & Raw gyro rate measurement \\
$\gamma$ & Complementary-filter weight (accelerometer term) \\
$\Delta t$ & Control/estimation loop sample period \\
$x_f, \alpha$ & Offset-correction filter state, filter coefficient \\
$d$ & Constant tilt-sensor offset \\
$\theta_{\text{th}}$ & Tilt threshold, jump-up $\to$ balancing switch \\
$K_d^{(i)}$ & Scheduled gain at linearization point $i$ (extension) \\
\bottomrule
\end{tabular}
\end{table}

\subsubsection{Gain Scheduling \& Mode Transitions}

The controller switches from the open-loop jump-up sequence to the LQR law
once $|\hat\theta_b| < \theta_{\text{th}}$, with hysteresis around
$\theta_{\text{th}}$ to prevent mode chattering. The baseline uses a single
fixed gain $K_d$ at the upright equilibrium. If recovery degrades near the
edges of the operating range, where the small-angle linearization weakens,
gain scheduling will be added: additional gains $K_d^{(i)}$ computed at
several linearization points, with the controller interpolating between them
based on $\hat\theta_b$, reusing the same discrete structure and offset
filter.

\subsection{Generalization to the 3D Cube}
\label{sec:3dmodel}

The dynamics, identification workflow, and control approach validated on the
2D model extend to the full three-dimensional cube, where three orthogonal
reaction wheels must stabilize the body about its edge and corner equilibria.

\subsubsection{Full-Body Equations of Motion}

With $R\in SO(3)$ the cube's orientation and $\omega_b\in\mathbb{R}^3$ its
body-frame angular velocity, the attitude kinematics are
\begin{equation}
\dot R = R\,[\omega_b]_\times,
\end{equation}
where $[\cdot]_\times$ is the skew-symmetric cross-product matrix. The
assembly's angular momentum about the pivot, expressed in the body frame,
generalizes the single-axis momentum term of the 2D model:
\begin{equation}
H = I\omega_b + \sum_{i=1}^{3} I_{w,i}\,\omega_{w,i}\,e_i.
\end{equation}
Euler's equation gives
\begin{equation}
\dot H + \omega_b \times H = \tau_g - C_b\,\omega_b,
\end{equation}
and each wheel obeys
\begin{equation}
I_{w,i}\big(\dot\omega_{b,i} + \dot\omega_{w,i}\big) = T_{m,i} - C_{w,i}\,\omega_{w,i}, \qquad i=1,2,3.
\end{equation}
Table~\ref{tab:3d_vars} defines all symbols.

\begin{table}[h]
\centering
\caption{Full-body (3D) model variables.}
\label{tab:3d_vars}
\begin{tabular}{ll}
\toprule
Symbol & Description \\
\midrule
$R$ & Cube orientation relative to world frame \\
$\omega_b$ & Body angular velocity, body frame ($\in\mathbb{R}^3$) \\
$\omega_{b,i} = e_i^{\mathrm T}\omega_b$ & Body angular velocity about wheel axis $i$ \\
$\omega_w = (\omega_{w,1},\omega_{w,2},\omega_{w,3})$ & Wheel angular velocities relative to body \\
$e_i$ & Spin axis of wheel $i$ (face normal) \\
$I$ & $3\times3$ inertia tensor, body+wheels, about pivot \\
$H$ & Total angular momentum about pivot \\
$\tau_g = m\,r_{cm}\times(R^{\mathrm T}\hat g)$ & Gravity torque about pivot \\
$m$ & Total assembly mass \\
$r_{cm}$ & Body-frame vector, pivot to assembly CoM \\
$\hat g$ & World vertical unit vector \\
$C_b$ & $3\times3$ body damping matrix \\
$I_{w,i}, C_{w,i}$ & Inertia, damping of wheel $i$ \\
$T_{m,i} = K_{m,i}u_i$ & Motor torque, wheel $i$ \\
$u_i, K_{m,i}$ & Motor current, torque constant, wheel $i$ \\
\bottomrule
\end{tabular}
\end{table}

Restricting motion to a single vertical plane with two wheel torques set to
zero recovers the 2D prototype model exactly, with $I_b, l_b, m_b, m_w, I_w,
C_b, C_w$ the corresponding single-axis projections of $I, r_{cm}, m, I_{w,i},
C_{w,i}$.


\subsubsection{Edge vs.\ Corner Equilibria \& Linearization}

Jump-up proceeds in two stages: Flat-to-Edge, where one wheel brakes,
followed by Edge-to-Corner, where a second wheel brakes. \textbf{Edge
balancing} constrains the body to a single axis, so the dynamics reduce to
the 2D form, and the 2D LQR baseline and offset filter apply directly, with
$m_b, l_b, I_b$ replaced by their edge equivalents.

\textbf{Corner balancing} is genuinely 3D: all three wheels contribute across
three rotational DOF. Linearizing about the corner-up equilibrium
$(\theta,\omega_b,\omega_w)=(0,0,0)$, with $\theta\in\mathbb{R}^3$ the small
orientation error, gives the same state-space form as the 2D case:
\begin{equation}
\dot x(t) = Ax(t) + Bu(t), \qquad x := (\theta,\omega_b,\omega_w) \in \mathbb{R}^9,
\end{equation}
with $A \in \mathbb{R}^{9\times9}$ and $B \in \mathbb{R}^{9\times3}$ built
from $I$, $r_{cm}$, and the per-wheel parameters $I_{w,i}, C_{w,i}, K_{m,i}$
(Table~\ref{tab:3d_vars}). The LQR design generalizes directly: $K_d$ becomes
a $3\times9$ gain matrix, using the same discretization, cost function, and
offset filter as Table~\ref{tab:control2d_vars}, with $Q, R$ weighting all
three tilt axes, body rates, and wheel speeds. The mode-transition logic
gains a third mode, Edge-to-Corner jump-up, triggered once edge balancing is
stable, followed by a second threshold switch into the corner LQR once
estimated orientation nears corner-up.